\documentclass[11pt,a4paper]{article}

% ── Packages ──────────────────────────────────────────────────────────────────
\usepackage[margin=1in]{geometry}
\usepackage{times}
\usepackage{microtype}
\usepackage{graphicx}
\usepackage{booktabs}
\usepackage{multirow}
\usepackage{amsmath}
\usepackage{amssymb}
\usepackage{array}
\usepackage{xcolor}
\usepackage{hyperref}
\usepackage{natbib}
\usepackage{caption}
\usepackage{subcaption}
\usepackage{float}
\usepackage{enumitem}
\usepackage{setspace}
\usepackage{fancyhdr}
\usepackage{abstract}
\usepackage{titlesec}

% ── Colors ────────────────────────────────────────────────────────────────────
\definecolor{arxivblue}{RGB}{30,80,155}
\definecolor{darkgray}{RGB}{60,60,60}

% ── Hyperref setup ────────────────────────────────────────────────────────────
\hypersetup{
  colorlinks=true,
  linkcolor=arxivblue,
  citecolor=arxivblue,
  urlcolor=arxivblue,
  pdftitle={Early Detection of At-Risk Students in Higher Education},
  pdfauthor={},
}

% ── Section formatting ────────────────────────────────────────────────────────
\titleformat{\section}{\large\bfseries\color{arxivblue}}{\thesection}{1em}{}
\titleformat{\subsection}{\normalsize\bfseries}{\thesubsection}{1em}{}
\titleformat{\subsubsection}{\normalsize\itshape}{\thesubsubsection}{1em}{}

% ── Header/footer ─────────────────────────────────────────────────────────────
\pagestyle{fancy}
\fancyhf{}
\fancyhead[L]{\small\textit{Early Detection of At-Risk Students Using Machine Learning}}
\fancyhead[R]{\small\thepage}
\renewcommand{\headrulewidth}{0.4pt}

% ── Caption style ─────────────────────────────────────────────────────────────
\captionsetup{font=small, labelfont=bf, skip=6pt}

% ── Column type for tables ────────────────────────────────────────────────────
\newcolumntype{C}[1]{>{\centering\arraybackslash}p{#1}}
\newcolumntype{L}[1]{>{\raggedright\arraybackslash}p{#1}}

% ─────────────────────────────────────────────────────────────────────────────
\begin{document}

% ── Title block ───────────────────────────────────────────────────────────────
\title{%
  \vspace{-0.5cm}
  \LARGE\bfseries Early Detection of At-Risk Students in Higher Education:\\
  \large A Comparative Study of Machine Learning Approaches for\\
  Binary and Multiclass Student Outcome Prediction
}

\author{%
  Gleb Babichev\\
  \small Tamos Education, Almaty, Kazakhstan\\
  \small \texttt{gleb.babichev007@gmail.com}
}

\date{March 2026 \quad|\quad arXiv preprint \quad|\quad
  \href{https://github.com/glebbabichevv/atrisk-students-ml}{\texttt{github.com/glebbabichevv/atrisk-students-ml}}}

\maketitle
\thispagestyle{fancy}

% ── Abstract ──────────────────────────────────────────────────────────────────
\begin{abstract}
\noindent
Student dropout represents one of the most pressing challenges in contemporary higher
education, carrying profound economic, social, and institutional consequences. This paper
presents a comprehensive empirical study on the early detection of at-risk students using
machine learning techniques applied to the publicly available UCI \textit{Predict Students'
Dropout and Academic Success} dataset ($N = 4{,}424$; \citealt{realinho2022}). We
systematically evaluate six classifiers---Logistic Regression, Decision Tree, Random Forest,
Support Vector Machine (SVM), Gradient Boosting, and Multilayer Perceptron
(MLP)---across three experimental settings: (1) \textit{binary classification} (Dropout
vs.\ Graduate, $n = 3{,}630$), (2) \textit{multiclass classification} (Dropout / Enrolled
/ Graduate, $n = 4{,}424$), and (3) \textit{early prediction} using only first-semester
features. Model performance is measured using accuracy, macro-averaged precision, recall,
F1-score, and AUC-ROC, with 95\% confidence intervals estimated via non-parametric bootstrap
resampling (1,000 iterations). In the binary setting, Logistic Regression achieves the
highest accuracy (0.9149) and AUC-ROC (0.9516), while Random Forest attains the best
macro F1-score (0.9046). In the multiclass setting, the presence of the \textit{Enrolled}
class---students with as-yet-undetermined outcomes---reduces macro F1 by approximately 20
percentage points across all models, with Gradient Boosting performing best ($\text{F1} =
0.7049$, $\text{AUC} = 0.8895$). Early prediction using first-semester data retains
87--92\% of full-model performance, demonstrating practical viability for timely
intervention. A fairness analysis reveals performance disparities across gender and
scholarship status subgroups. Our findings contribute a rigorous multi-setting benchmark
to the Educational Data Mining (EDM) literature and provide actionable guidance for
institutions deploying early warning systems.

\medskip\noindent
\textbf{Keywords:} student dropout prediction, machine learning, educational data mining,
early warning systems, binary vs.\ multiclass classification, gradient boosting,
bootstrap confidence intervals, learning analytics, algorithmic fairness.
\end{abstract}

\vspace{0.3cm}
\hrule
\vspace{0.3cm}

% ─────────────────────────────────────────────────────────────────────────────
\section{Introduction}
\label{sec:intro}

Student dropout in higher education is a global phenomenon with far-reaching consequences.
Institutions worldwide lose a significant proportion of enrolled students before degree
completion, resulting in wasted educational investment, reduced individual economic
potential, and diminished institutional reputation \citep{kemper2020}. According to
UNESCO data, dropout rates in higher education across OECD countries average between 30\%
and 50\%, with considerable variation across disciplines and institutional types. At the
individual level, dropout is associated with long-term reductions in lifetime earnings,
employment prospects, and social mobility \citep{chung2019}.

The emergence of Educational Data Mining (EDM) and Learning Analytics (LA) has provided
higher education institutions with powerful computational tools for understanding and
predicting student outcomes. \citet{romero2010} define EDM as the application of data
mining techniques to educational data to better understand students and the settings in
which they learn. Within this domain, dropout prediction has become one of the most active
research areas, with growing evidence that machine learning classifiers can identify at-risk
students with high accuracy, often before the end of the first semester
\citep{jayaprakash2014, macfadyen2010}.

Despite the proliferation of studies in this space, several methodological gaps remain.
First, many studies frame dropout prediction as a \textit{binary} classification
problem---contrasting eventual dropouts against graduates---while excluding students who
are still enrolled at the time of data collection. While this exclusion is
methodologically defensible (enrolled students have indeterminate outcomes), it raises
important questions about what happens when such students are retained in the dataset as a
third class. Second, few studies systematically compare binary and multiclass formulations
of the same prediction problem on the same dataset, making it difficult to understand the
performance cost of including unresolved cases. Third, many studies report point estimates
of accuracy or F1-score without confidence intervals, making it impossible to assess the
statistical reliability of performance differences between models \citep{efron1997}.

This paper addresses these gaps through a rigorous multi-setting benchmark study. We use
the UCI \textit{Predict Students' Dropout and Academic Success} dataset
\citep{realinho2022}, which contains records for 4,424 students enrolled at the
Polytechnic Institute of Portalegre, Portugal. The dataset covers 36 features spanning
demographic characteristics, socioeconomic background, academic history, and curricular
performance across two semesters. We evaluate six classifiers across three experimental
settings and report all metrics with 95\% bootstrap confidence intervals. A fairness
analysis examines model performance across gender and scholarship status subgroups.

\subsection*{Research Questions}

This study is guided by four research questions:
\begin{enumerate}[leftmargin=1.5em, itemsep=2pt, topsep=4pt]
  \item[\textbf{RQ1}] How do six established classifiers compare on binary dropout
    prediction (Dropout vs.\ Graduate)?
  \item[\textbf{RQ2}] How does model performance change when the Enrolled class is
    retained as a third outcome, and what does this reveal about the inherent
    predictability of student trajectories?
  \item[\textbf{RQ3}] Can meaningful prediction be achieved using only first-semester
    academic data, enabling earlier institutional intervention?
  \item[\textbf{RQ4}] Are there fairness concerns in model performance across gender
    and scholarship status subgroups?
\end{enumerate}

\subsection*{Contributions}

The main contributions of this paper are: (1) a systematic binary vs.\ multiclass
comparison on a standardized benchmark dataset with bootstrap CIs; (2) an empirical
quantification of the performance cost of including enrolled students in the outcome
space; (3) an early prediction experiment demonstrating practical viability using only
first-semester features; and (4) a subgroup fairness decomposition for the best-performing
multiclass model.

% ─────────────────────────────────────────────────────────────────────────────
\section{Related Work}
\label{sec:related}

\subsection{Educational Data Mining and Learning Analytics}

The field of EDM emerged in the early 2000s as a systematic application of data analysis
techniques to educational databases. \citet{romero2007} provided the first comprehensive
survey of EDM methods applied between 1995 and 2005, documenting classification,
clustering, association rule mining, and sequential pattern mining. Their subsequent survey
\citep{romero2010} expanded this analysis to 306 references and established a widely used
taxonomy of EDM tasks. A more recent update \citep{romero2020} documents the exponential
growth of the field into MOOCs, adaptive learning systems, and virtual reality
environments.

\citet{baker2014} define Learning Analytics as the measurement, collection, and analysis
of data about learners and their contexts for optimizing learning environments. 
\citet{papamitsiou2014} conducted a systematic review of 40 empirical LA studies,
identifying prediction and student modeling as dominant research themes.
\citet{namoun2021} synthesized 62 papers from 2010--2020, finding that supervised
classification---particularly decision trees, random forests, and neural
networks---dominates the student performance prediction literature.

\subsection{Dropout Prediction Using Machine Learning}

Machine learning-based dropout prediction has a well-established tradition.
\citet{macfadyen2010} provided an early proof of concept, demonstrating that a logistic
regression model built on LMS tracking data could correctly identify 81\% of failing
students. \citet{arnold2012} deployed one of the first large-scale early warning
systems---Course Signals at Purdue University---which reached over 23,000 students and
was associated with improved persistence rates.

Subsequent work expanded both feature sets and model sophistication. \citet{chung2019}
applied machine learning to national high school data in South Korea, achieving 85.4\%
accuracy with a random forest classifier. \citet{howard2018} compared eight prediction
methods at University College Dublin, finding that weeks 5--6 of the first semester
represent the optimal prediction window. \citet{akcapinar2019} demonstrated that $k$-NN
classifiers could achieve 89\% accuracy from LMS data as early as week three.

More recent work has focused on ensemble methods. \citet{niyogisubizo2022} proposed a
two-layer stacked generalization ensemble using Random Forest, XGBoost, and Gradient
Boosting as base learners, demonstrating consistent improvement over individual component
models. \citet{kemper2020} reviewed 42 dropout prediction studies, identifying systematic
issues including data leakage, lack of external validation, and overreliance on accuracy as
the sole evaluation metric.

\subsection{Studies Using the UCI Dropout Dataset}

The dataset used in this study was introduced by \citet{realinho2022}, who established
initial benchmarks for the three-class prediction problem. \citet{villar2024} conducted
a comprehensive comparative study on this dataset evaluating seven classifiers including
XGBoost, CatBoost, and LightGBM with SMOTE oversampling and Optuna hyperparameter
optimization, achieving AUC exceeding 0.93. \citet{martins2023} extended the analysis
through a phased prediction approach, demonstrating substantial accuracy improvements as
additional semester data becomes available. \citet{gu2026} applied recursive feature
elimination to reduce the feature set from 35 to 20 attributes, achieving AUC of 0.891
with Gradient Boosting while maintaining interpretability through SHAP analysis.

\subsection{Explainability and Fairness}

\citet{lundberg2017} introduced SHAP (SHapley Additive exPlanations), which has been
widely adopted in educational ML. \citet{nagy2024} applied SHAP alongside LIME and
partial dependence plots to a CatBoost dropout prediction model, producing individualized
risk explanations suitable for counselor use. \citet{krueger2023} demonstrated that
SHAP-based explanations can support personalized intervention planning with AUC-PR
reaching 89.5\%.

\citet{baker2022} provided a comprehensive review of algorithmic bias in education,
identifying race, gender, and socioeconomic status as the most common axes of demographic
disparity. \citet{gardner2019} introduced the ABROCA metric for quantifying predictive
fairness across subgroups, finding substantial gender disparities in several MOOC baseline
models. \citet{hu2020} proposed a fairness-constrained prediction framework for
identifying at-risk students.

\subsection{Class Imbalance}

\citet{wongvorachan2023} systematically compared undersampling, oversampling, and SMOTE
methods on the High School Longitudinal Study, finding that the optimal strategy depends
on imbalance degree. \citet{mduma2023} compared five resampling techniques, finding
SMOTE-ENN achieves the best overall performance. In this study, we use macro-averaged
metrics to ensure equal weight is given to all classes regardless of frequency, rather
than applying resampling procedures.

% ─────────────────────────────────────────────────────────────────────────────
\section{Dataset and Problem Formulation}
\label{sec:data}

\subsection{Dataset Description}

We use the \textit{Predict Students' Dropout and Academic Success} dataset
\citep{realinho2022}, available from the UCI Machine Learning Repository
(DOI:~10.24432/C5MC89) under Creative Commons Attribution 4.0 License. The dataset
contains records for 4,424 students enrolled in undergraduate programs at the Polytechnic
Institute of Portalegre, Portugal, spanning 13 different undergraduate degrees collected
between 2008 and 2019.

Three target outcomes are represented: \textit{Graduate} ($n = 2{,}209$; 49.9\%),
\textit{Dropout} ($n = 1{,}421$; 32.1\%), and \textit{Enrolled} ($n = 794$; 17.9\%) ---
the latter comprising students who were still active at the time of data collection. The
dataset contains \textbf{no missing values}.

The 36 features span four categories:
\begin{itemize}[itemsep=2pt, topsep=4pt]
  \item \textit{Demographic and socioeconomic:} age at enrollment, gender, nationality,
    parental education and occupation, scholarship status, displaced student indicator,
    debtor status, tuition fees up-to-date, international student status.
  \item \textit{Pre-enrollment academic:} admission grade, prior qualification grade and
    type, application mode and order.
  \item \textit{First-semester curricular:} units enrolled, approved, evaluated, and grade
    average.
  \item \textit{Second-semester curricular:} same set of indicators as first semester.
\end{itemize}

\begin{table}[H]
\centering
\caption{Dataset class distribution.}
\label{tab:dataset}
\begin{tabular}{lccc}
\toprule
\textbf{Class} & \textbf{N} & \textbf{Proportion} & \textbf{Description} \\
\midrule
Graduate  & 2,209 & 49.9\% & Completed degree \\
Dropout   & 1,421 & 32.1\% & Left before completion \\
Enrolled  &   794 & 17.9\% & Still active (outcome unknown) \\
\midrule
\textbf{Total} & \textbf{4,424} & \textbf{100\%} & 36 features, 0 missing \\
\bottomrule
\end{tabular}
\end{table}

\subsection{Experimental Settings}

We define three experimental settings:

\textbf{Experiment A --- Binary Classification.} We exclude 794 Enrolled students and
train classifiers to distinguish Dropout (label 0) from Graduate (label 1). This setting,
standard in the EDM literature \citep{mduma2019, kemper2020}, considers only students with
known terminal outcomes ($n = 3{,}630$).

\textbf{Experiment B --- Multiclass Classification.} We retain all 4,424 records and
train classifiers to predict among three classes: Dropout (0), Enrolled (1), Graduate (2).
This setting reflects real-world deployment conditions where institutions must act on
currently enrolled students whose trajectories are not yet resolved.

\textbf{Experiment C --- Early Prediction.} We retain the multiclass formulation but
remove all six second-semester features ($p = 30$), testing whether actionable predictions
can be obtained before second-semester performance data is available.

\subsection{Preprocessing and Split}

All experiments use stratified 80/20 train/test splitting with fixed random seed (seed
= 42) to ensure reproducibility. Stratification preserves class proportions across
splits. \texttt{StandardScaler} (zero mean, unit variance) is fitted exclusively on
training data and applied to test data to prevent leakage. Scaling is applied to Logistic
Regression, SVM, and MLP; tree-based models receive unscaled input.

% ─────────────────────────────────────────────────────────────────────────────
\section{Methodology}
\label{sec:method}

\subsection*{Code Availability}

All code, figures, and the \LaTeX\ source for this paper are publicly available at:\\
\url{https://github.com/glebbabichevv/atrisk-students-ml}

\subsection{Classifiers}

We evaluate six classifiers representing a diverse range of learning paradigms, consistent
with our research proposal and the comparative literature
\citep{villar2024, realinho2022}:

\begin{enumerate}[leftmargin=2em, itemsep=3pt, topsep=4pt]
  \item \textbf{Logistic Regression (LR):} Linear discriminative model with softmax
    transformation (solver = lbfgs, $C = 1.0$, max\_iter = 1,000). Provides an
    interpretable linear baseline \citep{kemper2020}.

  \item \textbf{Decision Tree (DT):} Non-parametric model with recursive binary splitting
    (max\_depth = 10, min\_samples\_leaf = 5). Interpretable single-model baseline and
    building block for ensemble methods \citep{romero2010}.

  \item \textbf{Random Forest (RF):} Ensemble of 200 trees trained on bootstrap samples
    with random feature subsets (n\_estimators = 200, min\_samples\_leaf = 2). Combines
    variance reduction through bagging with implicit feature selection
    \citep{niyogisubizo2022}.

  \item \textbf{Support Vector Machine (SVM):} Kernel-based classifier finding the
    maximum-margin hyperplane (kernel = RBF, $C = 10$, $\gamma = $ scale,
    probability = True). Competitive on tabular educational tasks
    \citep{howard2018}.

  \item \textbf{Gradient Boosting (GB):} Sequential additive ensemble fitting each
    successive tree to the residuals of the previous model (n\_estimators = 200,
    learning\_rate = 0.05, max\_depth = 5). A close relative of XGBoost
    \citep{asselman2023, hamim2022}.

  \item \textbf{Multilayer Perceptron (MLP):} Feedforward neural network with hidden
    layers (256, 128, 64), ReLU activation, Adam optimizer, early stopping, max 500
    epochs \citep{vo2026, cho2023}.
\end{enumerate}

\subsection{Evaluation Metrics}

We use \textbf{macro-averaged} metrics that assign equal weight to each class regardless
of frequency, following the recommendation of \citet{kemper2020} for imbalanced
educational datasets:

\begin{itemize}[itemsep=2pt, topsep=4pt]
  \item \textbf{Accuracy:} Proportion of correctly classified instances.
  \item \textbf{Macro Precision / Recall / F1:} Unweighted mean of per-class values.
    F1 is our primary metric for multi-class comparison.
  \item \textbf{AUC-ROC (OvR):} Macro-averaged area under the ROC curve using
    One-vs-Rest decomposition \citep{realinho2022}.
\end{itemize}

\subsection{Bootstrap Confidence Intervals}

To quantify statistical uncertainty, we compute non-parametric bootstrap confidence
intervals \citep{efron1997} for all five metrics. We draw $B = 1{,}000$ bootstrap samples
with replacement from the test set, recompute each metric, and report the 2.5th and 97.5th
percentiles as 95\% CI bounds. This method, recommended by \citet{kohavi1995} and
\citet{molinaro2005}, requires no normality assumption. Bootstrap samples with fewer
distinct classes than expected are discarded.

\subsection{Fairness Analysis}

We conduct a post-hoc fairness analysis for the best multiclass model, decomposing macro
F1 by two sensitive attributes: \textit{Gender} (Female vs.\ Male) and \textit{Scholarship
Holder} (Yes vs.\ No), selected based on documented relevance to educational equity
\citep{baker2022, gardner2019}. We report per-subgroup F1 scores on the held-out test set
without applying fairness constraints during training.

% ─────────────────────────────────────────────────────────────────────────────
\section{Results}
\label{sec:results}

\subsection{Experiment A: Binary Classification}

Table~\ref{tab:binary} presents binary classification results with 95\% bootstrap CIs.
All models achieve high performance, reflecting the strong discriminability between terminal
outcomes when students with unresolved trajectories are excluded.

\begin{table}[H]
\centering
\caption{Binary classification results (Dropout vs.\ Graduate, $n_\text{test} = 726$).
  Bootstrap mean [95\% CI].}
\label{tab:binary}
\resizebox{\textwidth}{!}{%
\begin{tabular}{lC{2.8cm}C{2.8cm}C{2.8cm}C{2.8cm}C{2.8cm}}
\toprule
\textbf{Model} & \textbf{Accuracy} & \textbf{Precision} & \textbf{Recall} & \textbf{F1} & \textbf{AUC-ROC} \\
\midrule
Logistic Regression & \textbf{0.9149} [0.894--0.935] & 0.9201 [0.900--0.941] & 0.9013 [0.879--0.925] & 0.9087 [0.888--0.931] & \textbf{0.9516} [0.934--0.968] \\
Random Forest       & 0.9122 [0.891--0.933] & \textbf{0.9249} [0.907--0.944] & 0.8935 [0.870--0.916] & \textbf{0.9046} [0.882--0.926] & 0.9554 [0.940--0.970] \\
Gradient Boosting   & 0.9037 [0.883--0.924] & 0.9105 [0.890--0.930] & 0.8877 [0.864--0.911] & 0.8962 [0.873--0.918] & 0.9547 [0.936--0.970] \\
Neural Network      & 0.9026 [0.880--0.924] & 0.9147 [0.894--0.934] & 0.8832 [0.858--0.908] & 0.8941 [0.870--0.917] & 0.9519 [0.936--0.967] \\
SVM                 & 0.8998 [0.879--0.920] & 0.9032 [0.882--0.925] & 0.8858 [0.863--0.909] & 0.8926 [0.870--0.914] & 0.9441 [0.923--0.962] \\
Decision Tree       & 0.8834 [0.861--0.906] & 0.8837 [0.861--0.906] & 0.8699 [0.845--0.893] & 0.8754 [0.851--0.898] & 0.8977 [0.869--0.924] \\
\bottomrule
\end{tabular}}
\end{table}

Logistic Regression achieves the highest accuracy (0.9149) and AUC-ROC (0.9516). This
result is consistent with findings in the EDM literature that simple linear models can
outperform complex classifiers when the decision boundary is approximately linear and
features are properly scaled \citep{kemper2020}. Random Forest attains the highest macro
F1 (0.9046), indicating superior balance between precision and recall across both classes.
Importantly, the 95\% CIs of the top four models overlap substantially, suggesting
differences among LR, RF, GB, and MLP are not statistically meaningful at this test set
size.

\subsection{Experiment B: Multiclass Classification}

Table~\ref{tab:multi} presents multiclass results. The inclusion of the Enrolled class
produces a substantial and consistent performance degradation relative to the binary
setting, with macro F1 declining by approximately 20 percentage points across all
classifiers.

\begin{table}[H]
\centering
\caption{Multiclass classification results (Dropout / Enrolled / Graduate,
  $n_\text{test} = 885$). Bootstrap mean [95\% CI].}
\label{tab:multi}
\resizebox{\textwidth}{!}{%
\begin{tabular}{lC{2.8cm}C{2.8cm}C{2.8cm}C{2.8cm}C{2.8cm}}
\toprule
\textbf{Model} & \textbf{Accuracy} & \textbf{Precision} & \textbf{Recall} & \textbf{F1} & \textbf{AUC-ROC} \\
\midrule
Gradient Boosting   & 0.7665 [0.739--0.797] & 0.7203 [0.685--0.757] & 0.6964 [0.664--0.731] & \textbf{0.7049} [0.671--0.740] & \textbf{0.8895} [0.872--0.908] \\
Random Forest       & \textbf{0.7712} [0.745--0.799] & \textbf{0.7314} [0.694--0.767] & 0.6816 [0.653--0.713] & 0.6936 [0.661--0.728] & 0.8893 [0.871--0.907] \\
Logistic Regression & 0.7675 [0.740--0.796] & 0.7054 [0.666--0.743] & 0.6741 [0.643--0.705] & 0.6807 [0.645--0.714] & 0.8772 [0.858--0.897] \\
Neural Network      & 0.7499 [0.721--0.779] & 0.6869 [0.648--0.723] & 0.6642 [0.633--0.695] & 0.6700 [0.635--0.703] & 0.8598 [0.840--0.880] \\
Decision Tree       & 0.7105 [0.684--0.740] & 0.6488 [0.615--0.682] & 0.6468 [0.614--0.681] & 0.6473 [0.613--0.681] & 0.8010 [0.774--0.828] \\
SVM                 & 0.7185 [0.688--0.749] & 0.6497 [0.616--0.684] & 0.6389 [0.606--0.670] & 0.6426 [0.610--0.675] & 0.8448 [0.824--0.865] \\
\bottomrule
\end{tabular}}
\end{table}

This degradation reflects a fundamental epistemological challenge: Enrolled students
occupy an intermediate feature space position that overlaps with both Dropout and Graduate
distributions, creating irreducible classification difficulty we term \textit{outcome
uncertainty bias}. Gradient Boosting achieves the best multiclass F1 (0.7049) and
AUC-ROC (0.8895). The reversal of model rankings between binary and multiclass
settings---with Gradient Boosting surpassing Logistic Regression---suggests the Enrolled
class discrimination task benefits from non-linear interaction modeling provided by
boosting.

\begin{figure}[H]
  \centering
  \includegraphics[width=\textwidth]{fig07_model_comparison.png}
  \caption{Model performance comparison with 95\% bootstrap confidence intervals
    across all five metrics (multiclass setting).}
  \label{fig:comparison}
\end{figure}

\begin{figure}[H]
  \centering
  \includegraphics[width=\textwidth]{fig08_roc_curves.png}
  \caption{ROC curves for all six classifiers using One-vs-Rest decomposition
    (multiclass setting). Panels show Dropout, Enrolled, and Graduate classes separately.}
  \label{fig:roc}
\end{figure}

\begin{figure}[H]
  \centering
  \includegraphics[width=\textwidth]{fig09_confusion_matrices.png}
  \caption{Row-normalized confusion matrices for all six classifiers (multiclass setting).
    All models struggle most with the Enrolled class, which is systematically
    misclassified as either Dropout or Graduate.}
  \label{fig:cm}
\end{figure}

Figure~\ref{fig:cm} reveals a consistent pattern: all models achieve high recall for
Graduate (the majority class) but struggle most with Enrolled students. This is consistent
with findings by \citet{martins2023}, who observed that enrolled students typically have
intermediate academic performance that does not unambiguously signal a future terminal
outcome.

\subsection{Binary vs. Multiclass Comparison}

\begin{figure}[H]
  \centering
  \includegraphics[width=\textwidth]{fig13_binary_vs_multiclass.png}
  \caption{Binary vs.\ multiclass performance comparison across all six models.
    Error bars show 95\% bootstrap CI. The performance gap is consistent across models
    and metrics.}
  \label{fig:bvmc}
\end{figure}

Figure~\ref{fig:bvmc} provides a direct visual comparison. The performance gap is
consistent across all models and all three plotted metrics. Averaged across models,
including the Enrolled class reduces macro F1 by 0.218 points (21.8\%), accuracy by 0.160
points (16.0\%), and AUC-ROC by 0.062 points (6.2\%). The smaller AUC penalty reflects
that AUC measures discriminability rather than calibration---models retain reasonable
risk-ranking ability even when absolute Enrolled classification is poor.

This finding has an important practical implication: institutions that need to classify
\textit{currently enrolled} students face a substantially harder prediction problem than
the idealized binary setting used in most published research. Practitioners should
calibrate their expectations accordingly when reading binary-only benchmark results.

\subsection{Experiment C: Early Prediction}

\begin{table}[H]
\centering
\caption{Early prediction results (30 features, first semester only). Final column shows
  F1 retained relative to full multiclass model.}
\label{tab:early}
\begin{tabular}{lcccccc}
\toprule
\textbf{Model} & \textbf{Accuracy} & \textbf{Precision} & \textbf{Recall} & \textbf{F1} & \textbf{AUC-ROC} & \textbf{\% of Full F1} \\
\midrule
Gradient Boosting   & 0.732 & 0.664 & 0.644 & 0.649 & 0.865 & 92.1\% \\
Logistic Regression & 0.734 & 0.648 & 0.616 & 0.611 & 0.853 & 89.7\% \\
Neural Network      & 0.729 & 0.648 & 0.615 & 0.616 & 0.853 & 91.8\% \\
Random Forest       & 0.726 & 0.639 & 0.614 & 0.610 & 0.863 & 87.9\% \\
SVM                 & 0.697 & 0.627 & 0.617 & 0.620 & 0.818 & 96.5\% \\
Decision Tree       & 0.685 & 0.608 & 0.603 & 0.604 & 0.777 & 93.3\% \\
\bottomrule
\end{tabular}
\end{table}

\begin{figure}[H]
  \centering
  \includegraphics[width=0.9\textwidth]{fig11_early_vs_full.png}
  \caption{Full model vs.\ early prediction F1-score comparison. All models lose no more
    than 12.1 percentage points of F1 when second-semester features are removed.}
  \label{fig:early}
\end{figure}

These results demonstrate that early prediction using first-semester data is practically
viable: all models lose no more than 12.1 percentage points of F1 (Table~\ref{tab:early}).
This finding is particularly significant for intervention planning: if an institution can
reliably identify high-risk students at the end of the first semester, it gains a full
semester of additional intervention time. \citet{jayaprakash2014} demonstrated in a
large multi-institution trial that early alerts, when accompanied by proactive advising,
can significantly improve retention rates. \citet{howard2018} identified weeks 5--6 of
the first semester as the optimal prediction window.

\subsection{Feature Importance}

\begin{figure}[H]
  \centering
  \includegraphics[width=\textwidth]{fig10_feature_importance.png}
  \caption{Top 20 feature importances from Random Forest (left) and Gradient Boosting
    (right). Second-semester academic performance dominates both rankings.}
  \label{fig:fi}
\end{figure}

Figure~\ref{fig:fi} shows that both ensemble models consistently rank second-semester
academic performance features highest: \textit{Curricular units 2nd semester (approved)}
and \textit{Curricular units 2nd semester (grade)} appear in the top three for both
models. First-semester performance and age at enrollment also rank consistently high.
Financial indicators---\textit{Tuition fees up to date} and \textit{Scholarship
holder}---appear in the top 20, confirming findings by \citet{realinho2022} and
\citet{fernandes2019} that financial factors predict dropout beyond academic performance.

\subsection{Fairness Analysis}

\begin{figure}[H]
  \centering
  \includegraphics[width=0.85\textwidth]{fig12_fairness.png}
  \caption{Fairness analysis: macro F1 by gender and scholarship status (Gradient
    Boosting, best multiclass model). Disparities are present across both sensitive
    attributes.}
  \label{fig:fairness}
\end{figure}

Figure~\ref{fig:fairness} reveals performance disparities across both sensitive
attributes. Students with scholarship status show different classification performance
than non-scholarship students, potentially because scholarship selection correlates with
features that are themselves predictive of graduation. Following \citet{baker2022} and
\citet{hu2020}, we caution that deploying a dropout prediction system without monitoring
for differential error rates could inadvertently concentrate institutional interventions
on student groups already over-scrutinized, while under-serving groups where the model
performs worst.

% ─────────────────────────────────────────────────────────────────────────────
\subsection{Feature Selection Analysis (Sub-RQ 3)}

To address Sub-RQ 3, we systematically evaluated three feature selection methods applied
to the multiclass setting with Gradient Boosting as the base classifier: Variance
Threshold, SelectKBest (ANOVA $F$-score), and Recursive Feature Elimination (RFE) with
Random Forest.

\begin{table}[H]
\centering
\caption{Impact of feature selection on Gradient Boosting multiclass performance.}
\label{tab:featsel}
\begin{tabular}{lcccc}
\toprule
\textbf{Method} & \textbf{Features} & \textbf{Accuracy} & \textbf{F1} & \textbf{AUC-ROC} \\
\midrule
Baseline (all features)   & 36 & 0.7063 & 0.7063 & 0.8903 \\
Variance Threshold        & 36 & 0.7063 & 0.7063 & 0.8903 \\
SelectKBest ($k=10$)      & 10 & 0.7480 & 0.6651 & 0.8546 \\
SelectKBest ($k=15$)      & 15 & 0.7559 & 0.6872 & 0.8806 \\
SelectKBest ($k=20$)      & 20 & 0.7627 & 0.7000 & 0.8795 \\
SelectKBest ($k=25$)      & 25 & 0.7571 & 0.6872 & 0.8829 \\
SelectKBest ($k=30$)      & 30 & 0.7582 & 0.6922 & 0.8876 \\
RFE ($n=10$)              & 10 & 0.7288 & 0.6461 & 0.8666 \\
RFE ($n=15$)              & 15 & 0.7627 & 0.6953 & 0.8845 \\
RFE ($n=20$)              & 20 & 0.7548 & 0.6835 & 0.8852 \\
RFE ($n=25$)              & 25 & 0.7537 & 0.6885 & 0.8897 \\
\bottomrule
\end{tabular}
\end{table}

\begin{figure}[H]
  \centering
  \includegraphics[width=\textwidth]{fig14_feature_selection.png}
  \caption{Left: macro F1 by feature selection method. Right: F1 vs.\ number of selected
    features for SelectKBest and RFE. Dashed line shows baseline (36 features).}
  \label{fig:featsel}
\end{figure}

\begin{figure}[H]
  \centering
  \includegraphics[width=0.75\textwidth]{fig15_selectkbest_features.png}
  \caption{Top 15 features ranked by ANOVA $F$-score (SelectKBest, $k=20$). Academic
    performance features dominate, with financial indicators also prominent.}
  \label{fig:skbfeats}
\end{figure}

These results directly address Sub-RQ 3. Variance Threshold removes no features (all
features have sufficient variance), confirming the dataset is already well-conditioned.
SelectKBest with $k=20$ retains 99.1\% of baseline F1 (0.7000 vs.\ 0.7063) using only
55.6\% of features, demonstrating strong feature redundancy. RFE with $n=15$ similarly
achieves 98.4\% of baseline F1 with 15 features. Both methods confirm that the top
predictive signal is concentrated in a small subset of academic performance and financial
features, consistent with the feature importance rankings in Section~\ref{sec:results}.

The ANOVA $F$-score rankings (Figure~\ref{fig:skbfeats}) reveal a three-tier structure:
(1) curricular performance features (2nd and 1st semester approved units and grades)
dominate with $F$-scores exceeding 500; (2) financial indicators (\textit{Tuition fees up
to date}, \textit{Scholarship holder}, \textit{Debtor}) form a secondary tier; and (3)
demographic factors (age, gender, application mode) contribute at a lower but non-trivial
level. This hierarchy is consistent with the theoretical model proposed by \citet{kemper2020}
and the SHAP analysis presented in Section~\ref{subsec:shap}.

\subsection{SHAP Explainability Analysis (Sub-RQ 2)}
\label{subsec:shap}

To address Sub-RQ 2, we applied SHAP (SHapley Additive exPlanations;
\citealt{lundberg2017}) to the best-performing multiclass model. Due to incompatibility
of SHAP's \texttt{TreeExplainer} with sklearn's \texttt{GradientBoostingClassifier} in
the multiclass setting, SHAP analysis was conducted using XGBoost
\citep{asselman2023} with identical hyperparameters (n\_estimators=200,
learning\_rate=0.05, max\_depth=5). XGBoost achieved comparable multiclass F1 of 0.711,
consistent with the sklearn Gradient Boosting baseline of 0.705.

\begin{figure}[H]
  \centering
  \includegraphics[width=0.75\textwidth]{fig17_shap_global_importance.png}
  \caption{Global SHAP feature importance (mean $|\text{SHAP}|$) averaged across all
    three classes. \textit{Curricular units 2nd sem (approved)} dominates with a mean
    $|\text{SHAP}|$ of 0.647, more than three times the second-ranked feature.}
  \label{fig:shapglobal}
\end{figure}

Figure~\ref{fig:shapglobal} presents global SHAP importance averaged across all classes.
\textit{Curricular units 2nd sem (approved)} is the dominant feature by a substantial
margin (mean $|\text{SHAP}| = 0.647$), followed by \textit{Curricular units 1st sem
(approved)} (0.217) and \textit{Tuition fees up to date} (0.182). This three-feature
cluster accounts for the majority of the model's predictive signal, with all remaining
features contributing at least three times less.

\begin{figure}[H]
  \centering
  \begin{subfigure}[b]{0.32\textwidth}
    \includegraphics[width=\textwidth]{fig16_shap_beeswarm_dropout.png}
    \caption{Dropout}
  \end{subfigure}
  \hfill
  \begin{subfigure}[b]{0.32\textwidth}
    \includegraphics[width=\textwidth]{fig16_shap_beeswarm_enrolled.png}
    \caption{Enrolled}
  \end{subfigure}
  \hfill
  \begin{subfigure}[b]{0.32\textwidth}
    \includegraphics[width=\textwidth]{fig16_shap_beeswarm_graduate.png}
    \caption{Graduate}
  \end{subfigure}
  \caption{SHAP beeswarm plots per class. Each dot represents one student. Red = high
    feature value, blue = low. Horizontal position shows impact direction on model output.}
  \label{fig:beeswarm}
\end{figure}

Figure~\ref{fig:beeswarm} reveals class-specific feature dynamics. For the
\textit{Dropout} class, low values of \textit{Curricular units 2nd sem (approved)}
(blue dots extending right) strongly increase dropout probability, while
\textit{Tuition fees up to date} shows a striking pattern: students who are
\textit{not} up to date on fees (blue) have dramatically higher dropout risk, confirming
that financial precarity is a direct dropout driver \citep{realinho2022}. For the
\textit{Enrolled} class, the dominant features are \textit{evaluations} rather than
\textit{approved units}, reflecting that enrolled students are still in the process of
completing assessments. For the \textit{Graduate} class, \textit{Scholarship holder}
enters the top features (global rank 8), with scholarship recipients showing higher
graduation probability.

\begin{figure}[H]
  \centering
  \begin{subfigure}[b]{0.32\textwidth}
    \includegraphics[width=\textwidth]{fig18_shap_waterfall_dropout.png}
    \caption{Dropout student}
  \end{subfigure}
  \hfill
  \begin{subfigure}[b]{0.32\textwidth}
    \includegraphics[width=\textwidth]{fig18_shap_waterfall_enrolled.png}
    \caption{Enrolled student}
  \end{subfigure}
  \hfill
  \begin{subfigure}[b]{0.32\textwidth}
    \includegraphics[width=\textwidth]{fig18_shap_waterfall_graduate.png}
    \caption{Graduate student}
  \end{subfigure}
  \caption{SHAP waterfall plots for three individual correctly classified students.
    Each bar shows the additive contribution of one feature to the final prediction,
    starting from the baseline $E[f(X)]$.}
  \label{fig:waterfall}
\end{figure}

Figure~\ref{fig:waterfall} provides individualized explanations for three correctly
classified students. The dropout student (f(x) = 2.078) is primarily flagged by a very
low \textit{Curricular units 2nd sem (approved)} value ($-1.464$, SHAP $= +1.67$),
representing the single largest individual feature contribution in our analysis. The
enrolled student (f(x) = 0.446) is characterized by older age at enrollment
(\textit{Age at enrollment} $= 0.888$, SHAP $= -0.17$) and moderate but incomplete
assessment activity. The graduate student (f(x) = 2.412) shows high 2nd semester
approved units ($0.517$, SHAP $= +0.77$) alongside a strong admission grade ($2.28$,
SHAP $= +0.16$). These individual-level explanations demonstrate the potential of
SHAP-based profiles for personalized intervention planning \citep{nagy2024,krueger2023}.

\section{Discussion}
\label{sec:discussion}

\subsection{The Enrolled Class Problem}

The most important finding of this study is the large and consistent performance
degradation caused by including Enrolled students in the outcome space. We propose that
this degradation has both a \textit{statistical} component (the minority Enrolled class
adds noise and increases classification difficulty) and an \textit{epistemological}
component (no historical pattern of features can perfectly predict a future outcome that
has not yet occurred). The epistemological component is fundamental and cannot be resolved
by algorithmic improvement---it reflects genuine uncertainty about human decision-making
and life circumstances that unfold over time.

Our finding parallels \citet{martins2023}, who showed that prediction accuracy for
enrolled students improves substantially as data from additional semesters becomes
available. This suggests the optimal architecture for a real-world early warning system
may be a \textit{phased prediction} approach: deploying a binary-style model for
first-year students (when intervention is most valuable) while transitioning to a
multiclass formulation as trajectories become clearer in later semesters.

\subsection{Model Selection Implications}

Our results reveal that model rankings are not stable across experimental settings.
Logistic Regression is the top performer in the binary setting but ranks third in the
multiclass setting, where Gradient Boosting leads. This instability underlines the
importance of selecting evaluation settings that match deployment context. An institution
deploying a model to predict outcomes for currently enrolled students should benchmark
models in the multiclass setting, not the cleaner binary formulation used in most
published studies.

The strong binary performance of Logistic Regression also has practical implications:
institutions with limited computational resources or strong interpretability requirements
may find a well-tuned logistic regression as effective as more complex ensembles, aligning
with observations in the broader tabular ML literature \citep{kemper2020}.

\subsection{Early Warning Viability}

The early prediction experiment demonstrates that first-semester data retains 87--92\% of
full-model performance. Given that second-semester data is typically unavailable until
4--6 months into the academic year, the practical intervention window is substantially
extended by deploying early prediction models. \citet{jayaprakash2014} demonstrated in
the Open Academic Analytics Initiative (OAAI) that early alerts accompanied by proactive
advising significantly improve retention; \citet{howard2018} identified the mid-semester
period as the optimal deployment window.

\subsection{Limitations}

This study has several limitations:
\begin{itemize}[itemsep=2pt, topsep=4pt]
  \item \textbf{Single institution:} The dataset was collected from one Portuguese
    polytechnic. Generalization requires external validation across different institutional
    contexts \citep{mduma2019}.
  \item \textbf{No hyperparameter optimization:} All models use manually specified
    hyperparameters. Systematic optimization (e.g., Optuna) would likely improve
    performance, particularly for ensemble methods \citep{villar2024}.
  \item \textbf{No resampling:} We did not apply SMOTE or other techniques; future work
    should evaluate whether resampling improves Enrolled-class recall
    \citep{mduma2023, masood2024}.
  \item \textbf{SHAP via XGBoost proxy:} SHAP analysis required substituting XGBoost
    for sklearn Gradient Boosting due to TreeExplainer incompatibility with multiclass
    sklearn GB. While XGBoost achieved comparable F1 (0.711 vs.\ 0.705), results may
    differ slightly from a native sklearn GB explanation.
\end{itemize}

% ─────────────────────────────────────────────────────────────────────────────
\section{Conclusion}
\label{sec:conclusion}

This paper presented a multi-setting benchmark study of machine learning approaches for
student dropout prediction using the UCI \textit{Predict Students' Dropout and Academic
Success} dataset \citep{realinho2022}. We evaluated six classifiers across binary
classification, multiclass classification, and early prediction settings, reporting all
metrics with 95\% bootstrap confidence intervals and complementing the analysis with
feature importance and fairness decomposition.

Our key findings are: (1) in the binary setting, all six classifiers achieve strong
performance ($\text{F1} > 0.875$, $\text{AUC} > 0.897$), with Logistic Regression and
Random Forest leading (RQ1); (2) including Enrolled students in a three-class formulation
consistently reduces macro F1 by $\sim$20 percentage points, reflecting an
epistemological limit on predicting unresolved trajectories (RQ2); (3) Gradient Boosting
is the most robust multiclass model ($\text{F1} = 0.705$, $\text{AUC} = 0.890$); (4)
early prediction using first-semester data retains 87--92\% of full-model performance,
supporting deployment before second-semester data is available (RQ3); (5) feature
selection analysis demonstrates that 15--20 features retain over 98\% of baseline F1,
with predictive signal concentrated in academic performance and financial indicators
(RQ3); (6) SHAP analysis identifies \textit{Curricular units 2nd sem (approved)} as the
dominant global predictor (mean $|\text{SHAP}| = 0.647$), with class-specific patterns
revealing that financial precarity drives dropout while scholarship status promotes
graduation (RQ2); and (7) fairness analysis reveals performance disparities by gender
and scholarship status that warrant monitoring in institutional deployment (RQ4).

From a practical standpoint, higher education institutions should: select evaluation
settings matching their deployment context; not assume clean binary benchmarks transfer to
real-world performance; treat the early prediction setting as operationally valuable; and
monitor demographic subgroup performance as a routine component of system evaluation.

Future work will extend this study in three directions: (1) SHAP-based explainability
\citep{lundberg2017, nagy2024} for individualized risk profiles; (2) fairness-constrained
training \citep{hu2020, kizilcec2022} to reduce demographic performance disparities; and
(3) phased prediction architectures \citep{martins2023} that dynamically update risk
scores as additional semester data becomes available.

% ─────────────────────────────────────────────────────────────────────────────
\bibliographystyle{plainnat}
\bibliography{references}

\end{document}
